\scanpage{248}
\mspage{194}

\[
\unsure{\Phi|\mathscr{D}}_\infty:\underline{\mathrm{Fib}}_W S\xrightarrow{\deleted{\unsure{\Phi}}}\underline{\mathrm{Perf}}_W S\subset\underline{\mathrm{Fib}}_W S
\]
et un morphisme fonctoriel
\[
i_X:X\xrightarrow{\deleted{\unsure{i_X}}}\Phi_\infty(X),\qquad
\tikz[baseline=(M.base)]\node[draw,ellipse,inner sep=3pt](M){$i_X\in W_S^f$};\ \text{si }X\in\underline{\mathrm{Fib}}_W S
\]

Ce foncteur permet alors, en l'absence de $W(4\,\mathrm b)$, de faire ce que le foncteur $C$ refusait \unsure{obstinément|absolument}. (Il est vrai que $C$ a l'avantage d'être défini sur $\mathrm{Cat}_{/S}$ tout entier, et d'y être \unsure{vraisemblablement|raisonnablement} $=$ fort utile, quand on a $W(4\mathrm b)$. Mais $\Phi_\infty$ est, lui aussi, défini sur $\mathrm{Cat}_{/S}$ \deleted{tout} entier, ainsi que $i_X$
\[
\begin{aligned}
\Phi_\infty:\mathrm{Cat}_{/S}&\longrightarrow\underline{\mathrm{Perf}}_W S\\
i_{\unsure{S|X}}:\mathrm{id}&\to\Phi_\infty\qquad\text{i.e. }
\begin{cases}
i_X:X\longrightarrow\Phi_\infty(X),\\
i_X\in W_S
\end{cases}
\end{aligned}
\]
avec le seul handicap (inévitable) que si \deleted{$X$ \unsure{n'est}} $X$ n'est pas dans $\underline{\mathrm{Fib}}_W S$, on ne peut dire mieux que $i_X\in W_S$, à l'exclusion de $i_X\in W_S^f$. \deleted{Pourtant, il suffit de} Mais pour $C(X)$ c'était \unsure{forcément|précis.} pareil. Donc on a l'impression que $\Phi_\infty$ (ou au choix, $\Psi_\infty$) bonnet blanc, blanc bonnet, joue, \deleted{\unsure{le}} en présence de $W(6\ \mathrm{bis})$, le même rôle que $C$ en présence de $W(4\mathrm b)$.
% slanted
\margin{\uline{En fait} $\Phi_\infty=C$, si $C$ est défini \unsure{par} des chemins infinis d'\unsure{1|un} seul côté \unsure{seulement} ; si \unsure{$C$ défini|je définis} $C$ \unsure{via|avec} les chemins \uline{finis} de \unsure{long.} \unsure{quelconque} \unsure{des} deux \unsure{côtés}, il prend en effet ses \unsure{val.} \unsure{ds} $\underline{\mathrm{Perf}}_W S$.}
De plus (en présence des deux axiomes à la fois) le foncteur $C$ a-t-il l'avantage esthétique de prendre ses valeurs dans $\underline{\mathrm{Perf}}_W S$, au lieu de $\underline{\unsure{\mathrm{P}|\mathrm{Pf}}}_W S$ seulement. Mais il s'impose de
% archive sheet number: 247
